\MCsolution{ex:curve-dm}
For a rational normalization component, stability gives $m_v\geq3$,
so $\deg T_{\mathbb P^1}(-B_v)=2-m_v<0$.
For an elliptic component, $m_v\geq1$ and the degree is $-m_v<0$.
For $h_v\geq2$, the degree $2-2h_v-m_v$ is already negative
even if $m_v=0$. An unmarked elliptic curve instead has trivial
tangent bundle and a nonzero translation vector field. A finite
group may have several elements: the hyperelliptic involution of
a smooth genus-two curve gives such an example. Triviality means
that only the identity remains, which is a stronger condition
than the one needed for a Deligne--Mumford stack.

\MCsolution{ex:curve-dimension}
There are two rational normalization components, each with three
special points, so both are stable. The genus is
$0+0+3-2+1=2$. Each component contributes
$3\cdot0-3+3=0$ parameters: a three-pointed rational curve has
no moduli. Each of the three nodes contributes one smoothing
parameter. The deformation space therefore has dimension $3$,
in agreement with $3g-3=3$.
